% GENERATED FILE. Edit canonical lessons, then run npm run generate:books.
% canonical-source-hash: fnv1a64:174d28b5481d66bd
% canonical-lessons: IT-C309-giovinezza, IT-C309-maggiorenne, IT-C309-minorenne, IT-C309-generazione, IT-C309-tramandare

\chapter{Di padre in figlio}
\label{ch:it-di-padre-in-figlio}
\hlchaptermodality{309}

\begin{chapteropening}
\textbf{By the end of this chapter:} \emph{I can talk about youth, coming of age and the generations that hand things down.}

Complete the last lesson of chapter 309: I can talk about youth, coming of age and the generations that hand things down.
\end{chapteropening}

\section[la giovinezza]{la giovinezza — youth (the time of life)}

\label{lesson:IT-C309-giovinezza}



\begin{quote}
Before the new one: say the Italian for the godmother, then the Italian for the confirmation.
\end{quote}

\subsection*{You'll want to know: la giovinezza}
\textbf{la giovinezza} — \textquotedblleft{}youth (the time of life)\textquotedblright{}.

\begin{quote}
Nella mia giovinezza ho vissuto al mare.
\end{quote}

\begin{quote}
\emph{In my youth I lived by the sea.}
\end{quote}

Built on \emph{giovane}, ‘young’.

\begin{cousinweb}
\emph{giovinezza} is built on \emph{giovine}, an old form of \emph{giovane} \textquotedblleft{}young\textquotedblright{}, from Latin \emph{iuvenis} (Lewis \& Short).
\end{cousinweb}

\begin{grammarlens}[title={What you've built}]
One more word for telling, explaining and describing at B1.
\end{grammarlens}

\subsection*{Your turn}
\begin{itemize}
\raggedright
  \item \emph{Say it:} \emph{la giovinezza}
  \item \emph{Say it:} \emph{la giovinezza}, once more
  \item \emph{Say it:} a sentence of your own with \emph{la giovinezza}
  \item \emph{Recall:} say the Italian for the godmother, then the Italian for the confirmation, then say \emph{la giovinezza} again
\end{itemize}

\begin{tcolorbox}[breakable,colback=teal!4,colframe=teal!35!black,title={Before you move on}]
What does la giovinezza mean? (\textquotedblleft{}Youth (the time of life)\textquotedblright{}.) Say it once more.
\end{tcolorbox}
\section[maggiorenne]{maggiorenne — of age, adult (18 or over)}

\label{lesson:IT-C309-maggiorenne}



\begin{quote}
Before the new one: say the Italian for the confirmation, then the Italian for youth.
\end{quote}

\subsection*{You'll want to know: maggiorenne}
\textbf{maggiorenne} — \textquotedblleft{}of age, adult (18 or over)\textquotedblright{}.

\begin{quote}
Mio figlio è diventato maggiorenne.
\end{quote}

\begin{quote}
\emph{My son has come of age.}
\end{quote}

\emph{La maggiore età} — the age of majority.

\begin{cousinweb}
\emph{maggiorenne} joins \emph{maggiore} \textquotedblleft{}greater\textquotedblright{}, from Latin \emph{maior}, and \emph{-enne} \textquotedblleft{}-year-old\textquotedblright{}, from \emph{annus} \textquotedblleft{}year\textquotedblright{} (Lewis \& Short).
\end{cousinweb}

\begin{grammarlens}[title={What you've built}]
One more word for telling, explaining and describing at B1.
\end{grammarlens}

\subsection*{Your turn}
\begin{itemize}
\raggedright
  \item \emph{Say it:} \emph{maggiorenne}
  \item \emph{Say it:} \emph{maggiorenne}, once more
  \item \emph{Say it:} a sentence of your own with \emph{maggiorenne}
  \item \emph{Recall:} say the Italian for the confirmation, then the Italian for youth, then say \emph{maggiorenne} again
\end{itemize}

\begin{tcolorbox}[breakable,colback=teal!4,colframe=teal!35!black,title={Before you move on}]
What does maggiorenne mean? (\textquotedblleft{}Of age, adult (18 or over)\textquotedblright{}.) Say it once more.
\end{tcolorbox}
\section[minorenne]{minorenne — underage; a minor}

\label{lesson:IT-C309-minorenne}



\begin{quote}
Before the new one: say the Italian for youth, then the Italian for of age.
\end{quote}

\subsection*{You'll want to know: minorenne}
\textbf{minorenne} — \textquotedblleft{}underage; a minor\textquotedblright{}.

\begin{quote}
I minorenni non possono entrare.
\end{quote}

\begin{quote}
\emph{Minors cannot come in.}
\end{quote}

Also a noun: \emph{un minorenne}.

\begin{cousinweb}
\emph{minorenne} joins \emph{minore} \textquotedblleft{}lesser\textquotedblright{}, from Latin \emph{minor}, and \emph{-enne}, from \emph{annus} \textquotedblleft{}year\textquotedblright{} (Lewis \& Short).
\end{cousinweb}

\begin{grammarlens}[title={What you've built}]
One more word for telling, explaining and describing at B1.
\end{grammarlens}

\subsection*{Your turn}
\begin{itemize}
\raggedright
  \item \emph{Say it:} \emph{minorenne}
  \item \emph{Say it:} \emph{minorenne}, once more
  \item \emph{Say it:} a sentence of your own with \emph{minorenne}
  \item \emph{Recall:} say the Italian for youth, then the Italian for of age, then say \emph{minorenne} again
\end{itemize}

\begin{tcolorbox}[breakable,colback=teal!4,colframe=teal!35!black,title={Before you move on}]
What does minorenne mean? (\textquotedblleft{}Underage; a minor\textquotedblright{}.) Say it once more.
\end{tcolorbox}
\section[la generazione]{la generazione — the generation}

\label{lesson:IT-C309-generazione}



\begin{quote}
Before the new one: say the Italian for of age, then the Italian for underage.
\end{quote}

\subsection*{You'll want to know: la generazione}
\textbf{la generazione} — \textquotedblleft{}the generation\textquotedblright{}.

\begin{quote}
Tre generazioni vivono nella stessa casa.
\end{quote}

\begin{quote}
\emph{Three generations live in the same house.}
\end{quote}

\emph{Di generazione in generazione} — from generation to generation.

\begin{cousinweb}
\emph{generazione} is a learned borrowing from Latin \emph{generātiō} \textquotedblleft{}a begetting\textquotedblright{}, from \emph{generāre} \textquotedblleft{}to beget\textquotedblright{}, built on \emph{genus} \textquotedblleft{}birth, kind\textquotedblright{} (Lewis \& Short).
\end{cousinweb}

\begin{grammarlens}[title={What you've built}]
One more word for telling, explaining and describing at B1.
\end{grammarlens}

\subsection*{Your turn}
\begin{itemize}
\raggedright
  \item \emph{Say it:} \emph{la generazione}
  \item \emph{Say it:} \emph{la generazione}, once more
  \item \emph{Say it:} a sentence of your own with \emph{la generazione}
  \item \emph{Recall:} say the Italian for of age, then the Italian for underage, then say \emph{la generazione} again
\end{itemize}

\begin{tcolorbox}[breakable,colback=teal!4,colframe=teal!35!black,title={Before you move on}]
What does la generazione mean? (\textquotedblleft{}The generation\textquotedblright{}.) Say it once more.
\end{tcolorbox}
\section[tramandare]{tramandare — to hand down (across generations)}

\label{lesson:IT-C309-tramandare}



\begin{quote}
Before the new one: say the Italian for underage, then the Italian for the generation.
\end{quote}

\subsection*{You'll want to know: tramandare}
\textbf{tramandare} — \textquotedblleft{}to hand down (across generations)\textquotedblright{}.

\begin{quote}
La nonna ci ha tramandato questa ricetta.
\end{quote}

\begin{quote}
\emph{Grandmother handed this recipe down to us.}
\end{quote}

For stories, recipes and customs.

\begin{cousinweb}
\emph{tramandare} is usually traced to a Latin \emph{trānsmandāre}, from \emph{trāns} \textquotedblleft{}across\textquotedblright{} and \emph{mandāre} \textquotedblleft{}to entrust\textquotedblright{} (Lewis \& Short gives \emph{mandāre}): to send across the years.
\end{cousinweb}

\begin{grammarlens}[title={What you've built}]
One more word for telling, explaining and describing at B1.
\end{grammarlens}

\subsection*{Your turn}
\begin{itemize}
\raggedright
  \item \emph{Say it:} \emph{tramandare}
  \item \emph{Say it:} \emph{tramandare}, once more
  \item \emph{Say it:} a sentence of your own with \emph{tramandare}
  \item \emph{Recall:} say the Italian for underage, then the Italian for the generation, then say \emph{tramandare} again
\end{itemize}

\begin{tcolorbox}[breakable,colback=teal!4,colframe=teal!35!black,title={Before you move on}]
What does tramandare mean? (\textquotedblleft{}To hand down (across generations)\textquotedblright{}.) Say it once more.
\end{tcolorbox}
